% filepath: /Users/sreeganeshbalasubramani/git/bhm/examples/toy_model/BHM_toy_explanation.tex
\documentclass[12pt]{article}

\usepackage{amsmath}
\usepackage{amssymb}
\usepackage{graphicx}
\usepackage{hyperref}
\usepackage{listings}
\usepackage{xcolor}
\usepackage{geometry}
\geometry{a4paper, margin=1in}

\lstset{
  basicstyle=\ttfamily\footnotesize,
  breaklines=true,
  frame=single,
  language=Python,
  keywordstyle=\color{blue},
  commentstyle=\color{gray},
  stringstyle=\color{orange}
}

\title{Detailed Explanation of \texttt{BHM\_toy.py}}
\author{Your Name}
\date{\today}

\begin{document}

\maketitle

\begin{abstract}
This document provides a detailed explanation of the Python script \texttt{BHM\_toy.py}. The script sets up a toy model framework for Bayesian Hierarchical Modeling (BHM), defines sets of coordinates in 3D space, calculates negative log-likelihoods for pairs of points, and proposes random moves for a Markov chain approach. Each function and significant code block is explained in detail.
\end{abstract}

\tableofcontents

\section{Introduction}
\texttt{BHM\_toy.py} is a Python script designed to illustrate or test hierarchical modeling concepts in a toy setting. The script:

\begin{itemize}
    \item Imports important libraries (NumPy, Matplotlib, SciPy, etc.).
    \item Defines arrays of points in 3D space for an idealized symmetrical system.
    \item Implements functions for calculating negative log-likelihood based on distances.
    \item Proposes moves in parameter space for potential sampling or optimization (e.g., in a Markov chain).
\end{itemize}

This document explains each part of the code in detail and describes how these parts fit together.

\section{Imports and Setup}
Below is a partial listing of the first lines of code:

\begin{lstlisting}
import numpy as np
import pickle
import matplotlib.pyplot as plt
from mpl_toolkits.mplot3d import Axes3D
from matplotlib.patches import Circle
import mpl_toolkits.mplot3d.art3d as art3d
from scipy.spatial.distance import cdist
from dataclasses import dataclass
from typing import Dict, List, Tuple
import random
\end{lstlisting}

\begin{itemize}
  \item \texttt{numpy} is used for numerical operations (e.g., array manipulations).
  \item \texttt{pickle} may be used for saving and loading Python objects.
  \item \texttt{matplotlib.pyplot} is for plotting in both 2D and 3D.
  \item \texttt{mpl\_toolkits.mplot3d.Axes3D} and \texttt{art3d} help with 3D rendering.
  \item \texttt{Circle} is a 2D patch that can be added to 3D plots for visual aids.
  \item \texttt{scipy.spatial.distance.cdist} computes pairwise distances between two sets of points.
  \item \texttt{dataclass} and \texttt{typing} assist in structuring data with type hints.
  \item \texttt{random} is used for generating random moves or seeds.
\end{itemize}

\section{Coordinate Arrays}
The script defines several NumPy arrays with 3D coordinates. For example:

\begin{lstlisting}
array_A = np.array([[0, 100, 0],
                   [-70.71, 70.71, 0],
                   ...
                   [70.71, 70.71, 0]])
\end{lstlisting}

\begin{itemize}
  \item These arrays (\texttt{array\_A}, \texttt{array\_B}, \texttt{array\_C}, etc.) specify a set of points in a 3D space at different \texttt{z}-heights or arranged in a geometric pattern.
  \item They appear to be arranged in a circular or symmetrical fashion (e.g., an 8-fold symmetry).
  \item Code references to these arrays are used later for distance calculations and for evaluating how well a configuration matches an “ideal” symmetrical arrangement.
\end{itemize}

\section{Core Class or Functions}

Although not all details might be shown in a single excerpt, the code typically defines a class or set of functions to handle:

\begin{itemize}
  \item \textbf{Positions:} A dictionary containing the current coordinates of all components or points.
  \item \textbf{Counts:} How many copies of each component type exist.
  \item \textbf{Distance Tolerances / Target Distances:} A dictionary specifying expected distances and tolerances between points.
  \item \textbf{Likelihood Computation:} Methods to calculate negative log-likelihoods based on how far each pair of points is from its expected distance.
  \item \textbf{Move Proposal:} A method to randomly perturb the positions to explore the parameter space, often used in MCMC.
\end{itemize}

\subsection{\texttt{calculate\_pair\_neg\_log\_likelihood\_matrix}}
In some part of the script (not fully shown in the excerpt), a function calculates a matrix of negative log-likelihoods for each pair of points:

\begin{lstlisting}
neg_log_likelihood_matrix = self.calculate_pair_neg_log_likelihood_matrix(
    positions[type1], 
    positions[type2],
    target_dist,
    tolerance
)
\end{lstlisting}

\begin{itemize}
  \item \texttt{positions[type1]} and \texttt{positions[type2]} are arrays storing the current 3D coordinates for each copy of a given component type.
  \item \texttt{target\_dist} is the expected distance between these two types of points.
  \item \texttt{tolerance} is how much deviation from \texttt{target\_dist} is allowed before penalizing the likelihood.
  \item This function likely computes pairwise distances (using \texttt{cdist}) and then transforms them into negative log-likelihood penalties.
\end{itemize}

\subsection{\texttt{Masking Symmetrical Off-Diagonal Terms}}
When \texttt{type1 == type2}, the code applies an upper-triangular mask to avoid double-counting pairs:

\begin{lstlisting}
if type1 == type2:
    mask = np.triu(np.ones_like(neg_log_likelihood_matrix), k=1)
    neg_log_likelihood_matrix *= mask
\end{lstlisting}

\begin{itemize}
    \item \texttt{np.triu(...)} creates a matrix with ones on/above the main diagonal (and zeros below).
    \item Multiplying by this mask ensures that only unique pairs are considered.
\end{itemize}

\subsection{\texttt{Unique Pairs and Summation}}
Next, the code finds the row minima and column minima of the negative log-likelihood matrix to identify the best pairing or alignment. It gathers these pairs in a set and sums the penalty:

\begin{lstlisting}
unique_pairs = set()
for i, j in enumerate(np.argmin(..., axis=1)):
    unique_pairs.add((i, j))
for i in np.argmin(..., axis=0):
    unique_pairs.add((i, j))

for i, j in unique_pairs:
    neg_log_likelihood += neg_log_likelihood_matrix[i, j]
\end{lstlisting}

\begin{itemize}
    \item \texttt{np.argmin(..., axis=1)} returns the index of the lowest value in each row.
    \item \texttt{np.argmin(..., axis=0)} does the same for each column.
    \item The set \texttt{unique\_pairs} ensures pairs aren’t duplicated when summing the negative log-likelihood.
\end{itemize}

\section{Proposed Moves}
The function \texttt{propose\_move} randomly perturbs positions:

\begin{lstlisting}
def propose_move(self, positions):
    new_positions = {k: v.copy() for k, v in positions.items()}

    # Select random component type and index
    type_name = np.random.choice(list(self.counts.keys()))
    idx = np.random.randint(self.counts[type_name])
\end{lstlisting}

\begin{itemize}
  \item \texttt{new\_positions} is a copy of the original positions dictionary to avoid modifying it in place.
  \item A random \texttt{type\_name} is chosen from the dictionary of component types.
  \item A random index \texttt{idx} is chosen, representing one specific copy of that component type.
  \item The chosen point's coordinates are then shifted by a random small amount (not shown in the excerpt), creating a candidate move for an MCMC or optimization algorithm.
\end{itemize}

\section{Negative Log-Likelihood}
Finally, the method computing the overall negative log-likelihood accumulates the penalty from every relevant pair. If the assembled structure is close to the target distances, the negative log-likelihood remains low. Otherwise, mismatches in the expected distances increase the penalty.

\section{Summary}
\begin{itemize}
  \item \textbf{Purpose:} \texttt{BHM\_toy.py} sets up a platform to explore hierarchical Bayesian modeling within a simplified or “toy” system of 3D points. 
  \item \textbf{Distance-Based Likelihood:} The code penalizes deviations from certain target distances through negative log-likelihood calculations. 
  \item \textbf{Random Moves:} A function proposes random moves of individual points, likely to be used in iterative or sampling-based algorithms (e.g., MCMC).
\end{itemize}

\noindent This toy setup is particularly useful for debugging and conceptual testing before moving on to more complex, real-world analysis.

% APPENDIX: Additional Explanation of BHM_toy.py (Lines 284-496)
% -------------------------------------------------------------------
% This LaTeX content can be added as a continuation to your existing
% "BHM_toy_explanation.tex" document.

\section{Detailed Explanation of the Code (Lines 284--496)}

\subsection{SystemParameters Dataclass}
\begin{lstlisting}[language=Python]
@dataclass
class SystemParameters:
    """Parameters for the molecular system"""
    box_size: float = 300.0
    radii: Dict[str, float] = None
    pair_distances: Dict[str, float] = None
    component_counts: Dict[str, int] = None
    
    def __post_init__(self):
        if self.radii is None:
            self.radii = {'A': 40.0, 'B': 10.0, 'C': 16.0}
        if self.pair_distances is None:
            self.pair_distances = {
                'AA': 80.5,
                'AB': 50.5,
                'CC': 31.5,
                'BC': 36.5
            }
        if self.component_counts is None:
            self.component_counts = {'A': 8, 'B': 8, 'C': 16}
\end{lstlisting}

\noindent
This \texttt{SystemParameters} dataclass holds parameters for the molecular system:
\begin{itemize}
    \item \textbf{\texttt{box\_size}:} The size of the simulation box in which components are placed (default: 300.0).
    \item \textbf{\texttt{radii}:} Radii (in \AA) of each component type (A, B, C). Default values are set if none are provided.
    \item \textbf{\texttt{pair\_distances}:} Target inter-component distances for pairs like AA, AB, CC, BC.
    \item \textbf{\texttt{component\_counts}:} How many copies of each component type exist in the system.
\end{itemize}
If the user does not pass these dictionaries, the \texttt{\_\_post\_init\_\_} method sets sensible defaults.

\subsection{BaseMCSampler Class}
\begin{lstlisting}[language=Python]
class BaseMCSampler:
    def __init__(self, params: SystemParameters):
        self.params = params
        self.positions = {}
        self.trajectory = []
        
        # Initialize sigma parameters with Jeffreys prior
        self.sigma = {}
        for pair_type in self.params.pair_distances.keys():
            # Sample from log-uniform distribution between 0.5 and 20 Å
            log_min, log_max = np.log(0.5), np.log(20.0)
            self.sigma[pair_type] = np.exp(
                np.random.uniform(log_min, log_max)
            )
            
        # Minimum likelihood value to avoid numerical issues
        self.min_likelihood = 1e-10
\end{lstlisting}

\begin{itemize}
    \item \textbf{\texttt{params}:} An instance of \texttt{SystemParameters}, storing system sizes, radii, etc.
    \item \textbf{\texttt{positions}:} A dictionary that will hold component positions (3D coordinates).
    \item \textbf{\texttt{trajectory}:} A list to store snapshots (states) for each Monte Carlo step.
    \item \textbf{\texttt{sigma}:} A dictionary of uncertainty (noise) parameters for each pair type (e.g., “AA”, “AB”). These are sampled initially from a log-uniform distribution between 0.5\AA{} and 20\AA{}, providing a broad prior.
    \item \textbf{\texttt{min\_likelihood}:} A small positive number (1e-10) to avoid taking the log of zero in likelihood calculations.
\end{itemize}

\subsubsection{\texttt{initialize\_positions}}
\begin{lstlisting}[language=Python]
def initialize_positions(self) -> Dict[str, np.ndarray]:
    """Initialize random positions for all components"""
    positions = {}
    for type_name, count in self.params.component_counts.items():
        positions[type_name] = np.random.uniform(
            0, self.params.box_size, (count, 3))
    return positions
\end{lstlisting}
Each component type is assigned an array of random coordinates, uniformly distributed in a 3D cube of size \texttt{box\_size}. This random initialization kickstarts the Monte Carlo procedure.

\subsubsection{\texttt{calculate\_exclusion\_score}}
\begin{lstlisting}[language=Python]
def calculate_exclusion_score(self, positions: Dict[str, np.ndarray]) -> float:
    """Calculate score term for excluded volume"""
    score = 0.0
    k_exclude = 100.0   # Strong exclusion penalty
    
    for type1 in positions:
        for type2 in positions:
            if type1 <= type2:
                pos1, pos2 = positions[type1], positions[type2]
                min_dist = self.params.radii[type1] + self.params.radii[type2]
                distances = cdist(pos1, pos2)
                ...
                if violations[0].size > 0:
                    score += k_exclude * \
                             np.sum((min_dist - distances[violations])**2)
    return score
\end{lstlisting}

\begin{itemize}
    \item The code computes a penalty if two particles overlap (i.e., the distance $< \texttt{min\_dist}$, which is the sum of their radii).
    \item \texttt{cdist(pos1, pos2)} calculates the pairwise distances between all copies of \texttt{type1} and all copies of \texttt{type2}.
    \item If \texttt{type1 == type2}, only the upper triangle of the distance matrix is considered to avoid double-counting.
    \item The penalty accumulates quadratically as $(\texttt{min\_dist} - \texttt{distance})^2$ under violations.
\end{itemize}

\subsubsection{\texttt{calculate\_pair\_score} and \texttt{calculate\_pair\_scores\_matrix}}
\begin{lstlisting}[language=Python]
def calculate_pair_score(self, pos1: np.ndarray, pos2: np.ndarray, 
                         target_dist: float, sigma: float) -> float:
    distance = np.sqrt(np.sum((pos1 - pos2)**2))
    score = ((distance - target_dist)**2) / (2 * sigma**2)
    return score

def calculate_pair_scores_matrix(self, pos1: np.ndarray, pos2: np.ndarray, 
                                 target_dist: float, sigma: float) 
                                 -> np.ndarray:
    distances = cdist(pos1, pos2)
    scores = ((distances - target_dist)**2) / (2 * sigma**2)
    return scores
\end{lstlisting}

\begin{itemize}
    \item The pair score (or “distance penalty”) is calculated as a Gaussian penalty around the \texttt{target\_dist}, scaled by $\sigma$.
    \item \texttt{calculate\_pair\_scores\_matrix} does the same for all pairs in two arrays of positions via \texttt{cdist}.
\end{itemize}

\subsubsection{\texttt{propose\_sigma\_move}}
\begin{lstlisting}[language=Python]
def propose_sigma_move(self) -> Dict[str, float]:
    """Propose new sigma values"""
    new_sigma = self.sigma.copy()
    # Pick random pair type
    pair_type = np.random.choice(list(self.sigma.keys()))
    
    # Propose new value in log space
    log_sigma = np.log(new_sigma[pair_type])
    step_size = 0.1
    new_log_sigma = log_sigma + np.random.normal(0, step_size)
    
    # Ensure within bounds [0.5, 20]
    new_log_sigma = np.clip(new_log_sigma, np.log(0.5), np.log(20.0))
    new_sigma[pair_type] = np.exp(new_log_sigma)
    
    return new_sigma
\end{lstlisting}

\begin{itemize}
    \item This method randomly adjusts one of the $\sigma$ values by sampling in log space.
    \item After shifting by a normal random variable of standard deviation \texttt{0.1}, the new $\sigma$ is clipped to [0.5, 20].
\end{itemize}

\subsubsection{\texttt{propose\_position\_move}}
\begin{lstlisting}[language=Python]
def propose_position_move(self, positions: Dict[str, np.ndarray]) -> Dict[str, np.ndarray]:
    """Propose new positions with adaptive step size"""
    new_positions = {k: v.copy() for k, v in positions.items()}
    
    # Select random component type and index
    type_name = np.random.choice(list(self.params.component_counts.keys()))
    idx = np.random.randint(self.params.component_counts[type_name])
    
    # Adaptive step size
    step_size = self.params.radii[type_name] * 0.1
    
    # Apply random displacement
    new_positions[type_name][idx] += np.random.normal(0, step_size, 3)
    
    # Ensure within box
    new_positions[type_name][idx] = np.clip(
        new_positions[type_name][idx], 0, self.params.box_size
    )
    
    return new_positions
\end{lstlisting}

\begin{itemize}
    \item A random component type and a random copy (index) of that type are chosen.
    \item The position is perturbed with a step size proportional to the component's radius.
    \item The new position is clipped so it remains within the simulation box (\texttt{[0, box\_size]}).
\end{itemize}

\subsubsection{\texttt{save\_state} and \texttt{save\_trajectory}}
\begin{lstlisting}[language=Python]
def save_state(self, step: int, positions: Dict[str, np.ndarray], 
               score: float) -> Dict:
    """Save current state information"""
    state = {
        'step': step,
        'positions': positions,
        'sigma': self.sigma.copy(),
        'score': score,
        'types': {},
        'bead_numbers': {}
    }
    # record type, bead index
    current_idx = 0
    for type_name, count in self.params.component_counts.items():
        for i in range(count):
            state['types'][current_idx] = type_name
            state['bead_numbers'][current_idx] = current_idx + 1
            current_idx += 1
    return state

def save_trajectory(self, filename: str = None):
    """Save trajectory to file"""
    if filename is None:
        filename = f"trajectory.pkl"
    with open(filename, 'wb') as f:
        pickle.dump(self.trajectory, f)
    return filename
\end{lstlisting}

\begin{itemize}
    \item \texttt{save\_state} collects information about the current Monte Carlo step: positions, $\sigma$ values, the step index, total score, etc.
    \item \texttt{save\_trajectory} writes the entire \texttt{trajectory} list to a pickle file for later analysis or visualization.
\end{itemize}

\subsubsection{\texttt{load\_trajectory}}
\begin{lstlisting}[language=Python]
@staticmethod
def load_trajectory(filename: str) -> List[Dict]:
    """Load trajectory from file"""
    with open(filename, 'rb') as f:
        return pickle.load(f)
\end{lstlisting}
\noindent
Loads the pickle file containing a trajectory list of states, returning it as a Python object.

\subsubsection{\texttt{run\_mc}}
\begin{lstlisting}[language=Python]
def run_mc(self, n_steps: int = 50000, save_freq: int = 100) -> Tuple:
    """Base Monte Carlo sampling with position and sigma moves"""
    positions = self.initialize_positions()
    current_score = self.calculate_score(positions)
    
    best_positions = {k: v.copy() for k, v in positions.items()}
    best_score = current_score
    
    self.trajectory = []
    accepted = 0
    
    # Temperature schedule
    initial_temp = 10.0
    final_temp = 0.1
    
    for step in range(n_steps):
        temp = initial_temp * \
               (final_temp / initial_temp)**(step / n_steps)
        
        # 90% chance position move, 10% sigma move
        if np.random.random() < 0.9:
            proposed_positions = self.propose_position_move(positions)
            proposed_sigma = self.sigma.copy()
        else:
            proposed_positions = positions.copy()
            proposed_sigma = self.propose_sigma_move()
        
        # Store old sigma, set new sigma
        old_sigma = self.sigma.copy()
        self.sigma = proposed_sigma
        
        # Calculate new score
        proposed_score = self.calculate_score(proposed_positions)
        delta_e = proposed_score - current_score
        
        # Metropolis criterion
        if delta_e < 0 or np.random.random() < np.exp(-delta_e/temp):
            positions = proposed_positions
            current_score = proposed_score
            accepted += 1
            
            if current_score < best_score:
                best_score = current_score
                best_positions = {k: v.copy() for k, v in positions.items()}
        else:
            self.sigma = old_sigma
        
        # Save trajectory info periodically
        if step % save_freq == 0:
            self.trajectory.append(self.save_state(step, positions, current_score))
            print(f"Step {step}, Score: {current_score:.2f}, "
                  f"T: {temp:.4f}, Accept: {accepted/(step+1):.2f}")
    
    return best_positions, self.trajectory, self.save_trajectory()
\end{lstlisting}

\begin{itemize}
    \item \textbf{Initialization:} Positions are set randomly, and a score is calculated.
    \item \textbf{Temperature Schedule:} A linear, geometric, or exponential cooling schedule can be used. Here, the code transitions from \texttt{initial\_temp}=10.0 to \texttt{final\_temp}=0.1 over \texttt{n\_steps}.
    \item \textbf{Move Selection:}
        \begin{itemize}
            \item With 90\% probability, perturb the \texttt{positions}.
            \item Otherwise, propose a new $\sigma$ for one pair type.
        \end{itemize}
    \item \textbf{Metropolis Criterion:} 
    \[
        \Delta E = \texttt{proposed\_score} - \texttt{current\_score}
        \quad \text{and} \quad
        p_{\text{accept}} = \exp \left(-\frac{\Delta E}{T}\right).
    \]
    If $\Delta E < 0$ (improvement) or if a random number $< p_{\text{accept}}$, the move is accepted.
    \item \textbf{Saving Trajectory:} Every \texttt{save\_freq} steps, the state is stored and a progress message is printed.
    \item \textbf{Return Value:} The function returns the best positions found, the full trajectory, and a call to \texttt{save\_trajectory()} that writes the trajectory to a file.
\end{itemize}

\subsection{Overall Workflow Summary}
\begin{enumerate}
    \item \textbf{Initialize} random coordinates for each component type within a cubic box.
    \item \textbf{Compute} an initial score (\texttt{calculate\_score} typically sums distance-based penalties and excluded-volume penalties).
    \item \textbf{Set Up} a Monte Carlo loop for a specified number of steps (\texttt{n\_steps}).
    \item \textbf{Propose Moves}:
    \begin{itemize}
        \item Position moves: Randomly pick a single component, perturb its coordinates, then evaluate the new score.
        \item Sigma moves: Randomly pick a single pair type and adjust its $\sigma$ parameter in log space.
    \end{itemize}
    \item \textbf{Accept or Reject Moves} using the Metropolis criterion, updating positions and $\sigma$.
    \item \textbf{Periodically Save} the current state (\texttt{positions}, \texttt{sigma}, \texttt{score}), log progress, and reduce temperature by a factor determined by the step fraction.
    \item \textbf{Return} the best positions discovered, the entire trajectory, and a saved file (pickle) of the trajectory for later analysis.
\end{enumerate}

\noindent
This concludes the detailed explanation of lines 284--496 in \texttt{BHM\_toy.py}. The code establishes a flexible Monte Carlo framework, allowing both positional shifts and parameter adjustments to seek configurations with minimal penalty (or maximal likelihood) under the model’s distance constraints and radii exclusions.

% End of Appendix

\end{document}